%\section{Multi-Party NIKE based on Zero-k-Clique with (Full) Security}
%\label{sec:hcp}
%In this section, we extend our multi-party NIKE to one with full security via privacy amplification. 
%We first introduce the probability bound as follows.
\begin{theorem}[Chebyshev bound]
%<<<<<<< HEAD
%<<<<<<< HEAD
%	Fix $\epsilon > 0$ and $b \in \{0,1\}$, and let $\{x_i\}$ be pairwise-independent, $0/1$-random variables for which $\Pr[x_i = b] \geq \frac{1}{2} + \epsilon$ for all $i$. 
%%	Consider the process in which $m$ values $x_1, \cdots, x_m$ are recorded and 
%	Let $x = \begin{cases} 1 & \text{if $\Sigma_{i=1}^{q}x_i \geq \frac{q}{2}$}\\
%	0 &  \text{if $\Sigma_{i=1}^{q}x_i < \frac{q}{2}$}
%	\end{cases}$. Then
%=======
	For any $\epsilon > 0$, $b \in \{0,1\}$, and $q\in\N$, let $\{x_i\}_{i\in[q]}$ be pairwise-independent, $0/1$-random variables such that
	$$\Pr[x_i = b] \geq 1/2 + \epsilon$$ for all $i\in[q]$.
%	Consider the process in which $m$ values $x_1, \cdots, x_m$ are recorded and $x$ is set to the value that occurs a strict majority of the time. Then
	Let $x = \begin{cases} 1 & \text{if $\Sigma_{i=1}^{q}x_i \geq q/2$}\\
	0 &  \text{if $\Sigma_{i=1}^{q}x_i < q/2$}
	\end{cases}$. Then
%>>>>>>> 460f90c11079c66ad36db4d2223f8b5c29dee719
	$$ \Pr[x \neq b] \leq 1/(4\epsilon^2q).$$
%=======
%	For any $\epsilon > 0$, $b \in \{0,1\}$ and $q\in\N$, and let $\{x_i\}_{i\in[q]}$ be pairwise-independent, $0/1$-random variables such that $\Pr[x_i = b] \geq \frac{1}{2} + \epsilon$ for all $i\in[q]$. Consider the process in which $m$ values $x_1, \cdots, x_m$ are recorded and $x$ is set to the value that occurs a strict majority of the time. Then
%	$$ \Pr[x \neq b] \leq \frac{1}{4\epsilon^2m}. $$
%>>>>>>> 00bd60cea94da8dde24d82e43fc4b4710637bcef
\end{theorem}

%appendix hcp theorem
%Let $\cc_1$ (respectively, $\cc_2$) be the set of all word-RAM families such that for each $\mathcal{F} = \{f_{\lambda}\}\in\cc_1$ and each $\lambda\in\N$, $f_{\lambda}$ runs in time $\tilde{O}(T_1(\lambda))$ (respectively, $\tilde{O}(T_2(\lambda))$) for any constant $\pnum \geq 2, k \geq 3$. Let $\epsilon =1/ \lambda^{o(1)}$ and $\tnike = (\gen, \share)$ be a $\cc_1$-$\pnum$-NIKE that satisfies $\cc_2$-$\frac{\epsilon}{4}$-weak security with the secret key size $ \log(\lambda^{\Omega(1)})$, we construct $\tnike^* = (\gen^*, \share^*)$ satisfying $\cc_2$-$\epsilon$-security as in \cref{fig:con-search-nike}.
%not \epsilon/4
%Let $\cc_1$ (respectively, $\cc_2$) be the set of all word-RAM families such that for each $\mathcal{F} = \{f_{\lambda}\}\in\cc_1$ and each $\lambda\in\N$, $f_{\lambda}$ runs in time $\tilde{O}(T_1(\lambda))$ (respectively, $\tilde{O}(T_2(\lambda))$) for any constant $\pnum \geq 2, k \geq 3$. Let $\tnike = (\gen, \share)$ be a $\cc_1$-$\pnum$-NIKE with shared-key length $Q(\lambda) = \log(\lambda^{\Omega(1)})$, and let $\tnike^* = (\gen^*, \share^*)$ be the construction in \cref{fig:con-search-nike}.
Let $\cc_1$ (respectively, $\cc_2$) be the set of all word-RAM families such that for each $\mathcal{F} = \{f_{\lambda}\}\in\cc_1$ and each $\lambda\in\N$, $f_{\lambda}$ runs in time $\tilde{O}(T_1(\lambda))$ (respectively, $\tilde{O}(T_2(\lambda))$) for any constant $\pnum \geq 2, k \geq 3$. Let $\epsilon =1/ \lambda^{o(1)}$ and $\tnike = (\gen, \share)$ be a $\cc_1$-$\pnum$-NIKE that satisfies $\cc_2$-$\epsilon$-weak security with the shared key size $ Q(\lambda) = \Theta(\log(\lambda))$, we construct $\tnike^* = (\gen^*, \share^*)$ satisfying $\cc_2$-$1/ \lambda^{o(1)}$-security as in \cref{fig:con-search-nike}.
%\begin{theorem}
%\label{thm:hcp}
%	Let $p_{\mathrm{GL}}(\epsilon)=\epsilon^{O(1)}$ denote the reconstruction success probability guaranteed by the standard Goldreich--Levin reconstruction algorithm for recovering a $Q(\lambda)$-bit string from an $\epsilon$-advantage predictor of its random inner products, using $\mathrm{poly}(Q(\lambda),1/\epsilon)$ oracle calls. Assume that $\mathrm{poly}(Q(\lambda),1/\epsilon)=\lambda^{o(1)}$. If $\tnike$ is a $\cc_1$-$\pnum$-NIKE with $\cc_2$-$\epsilon_w$-weak security and $\epsilon_w<p_{\mathrm{GL}}(\epsilon)$, then $\tnike^*$ is a $\cc_1$-$\pnum$-NIKE with $\cc_2$-$\epsilon$-security.
%\end{theorem}
\begin{theorem}
\label{thm:hcp}
	If $\tnike$ is a $\cc_1$-$\pnum$-NIKE with $\cc_2$-$\epsilon$-weak security, then $\tnike^*$ is a $\cc_1$-$\pnum$-NIKE with $\cc_2$-$1/ \lambda^{o(1)}$-security.
\end{theorem}
\begin{figure}[h!]
\begin{center}
\framebox{
\small
  \begin{tabular}{l|l}
    \begin{minipage}[t]{0.4\textwidth}
    	\underline{$\gen^*(\bot)$:}
	\item $(\pk, \sk) \getsr \gen(\bot)$
%	\item $\vec{v} \getsr \{0,1\}^{\log\elllambda^3}$
	\item $\vec{v} \getsr \{0,1\}^{Q(\lambda)}$
    	\item Return $(\pk^* =(\pk, \vec{v}),\sk^* = \sk)$
      \end{minipage}&  \begin{minipage}[t]{0.4\textwidth}
        
    \underline{$\share^*(i,\sk_i^*,(\pk_j^*)_{j\in[\pnum]})$:}
    \item Parse $\pk_j^*=(\pk_j,\vec v_j)$ for all $j \in [\pnum]$
    	\item $\tempkey = \share(i,\sk_i^*,(\pk_j)_{j\in[\pnum]})$
	%\item Let $\vec{v} = \vec{v}_\pnum$
	\item If $\tempkey=\bot$, return $\bot$
	\item Return $\key^* = \tempkey \cdot \vec{v}_\pnum$
    \end{minipage} 
      \end{tabular}
      }
\end{center}
\caption{Definition of $\tnike^*=\{\gen^*,\share^*\}$ with $\cc_2$-$1/ \lambda^{o(1)}$-security.}
\label{fig:con-search-nike}
\end{figure}
%\begin{proof}
%	Suppose there exists a distinguisher $\advA=\{\adv\}\in\cc_2$ that breaks the $\cc_2$-$\epsilon$-security of $\tnike^*$. Using $\adv$ as an oracle and varying the public vector $\vec v_\pnum$ in \cref{fig:con-search-nike}, the standard Goldreich--Levin reconstruction argument yields a randomized predictor that, on input the public keys of $\tnike$, recovers
%	\[
%	\key=\share(1,\sk_1,(\pk_i)_{i\in[\pnum]})
%	\]
%	with probability at least $p_{\mathrm{GL}}(\epsilon)$.
%	The reconstruction overhead is $\mathrm{poly}(Q(\lambda),1/\epsilon)$. By assumption this overhead is $\lambda^{o(1)}$, and is therefore absorbed by the $\widetilde{O}(\cdot)$ adversarial running time defining $\cc_2$. Hence, if $p_{\mathrm{GL}}(\epsilon)>\epsilon_w$, this predictor contradicts the $\cc_2$-$\epsilon_w$-weak security of $\tnike$. Therefore, $\tnike^*$ is $\cc_2$-$\epsilon$-secure.
%\end{proof}
\begin{proof}
	Let $m =  2Q(\lambda)/\epsilon^2$, $(\pk_i, \sk_i) \getsr \gen(\bot)$ for all $i \in [\pnum]$, $\key = \share(1, \sk_1, \pk_{[\pnum]})$, $\vec v_{\pnum} \getsr \bin^{Q(\lambda)}$, and $\advA = \{\adv\} \in \cc_2$ be any adversary breaking the $\cc_2$-$\epsilon$-security of $\tnike^*$.
%	where $\cc_2$ runs in time $\tilde{\Omega}(\elllambda T(\lambda))$. 
%	We construct an adversary $\advB = \{b_{\lambda}\} \in \cc_2$ breaking the $\cc_2$-$\frac{\epsilon}{4}$-weak security of $\tnike$.
	We construct an adversary $\advB = \{b_{\lambda}\} \in \cc_2$ breaking the $\cc_2$-$\epsilon$-weak security of $\tnike$.
	
	Firstly, $b_{\lambda}$ chooses $\log (m)$ pairs $(\sigma_1, r_1), \cdots, (\sigma_{\log (m)}, r_{\log (m)}) \getsr \{0,1\} \times \{0,1\}^{Q(\lambda)}$. Then, for every nonempty subset $\mathcal{X} \subseteq [\log (m)]$, $b_{\lambda}$ computes $r_\mathcal{X} = \oplus_{i \in \mathcal{X}} r_i$ and $\sigma_\mathcal{X} = \oplus_{i \in \mathcal{X}} \sigma_i$. For every $i \in [Q(\lambda)]$, $b_{\lambda}$ runs 
	$$\tempkey_i^\mathcal{X} = \sigma_\mathcal{X} \oplus \adv(\{\pk_i\}_{i\in[\pnum]}, r_\mathcal{X} \oplus e^{i}),$$ 
	where $e^{i}$ denotes the $Q(\lambda)$-bit string with $1$ in the $i$'th position and 0 everywhere else, and sets $\tempkey_i^*$ as the majority of ${\tempkey_i^\mathcal{X}}$. Finally, $b_{\lambda}$ outputs $\tempkey^* = \tempkey_1^*|| \cdots ||\tempkey_{Q(\lambda)}^*$.
	Let 
	$$\set = \{\tempkey \big{|} \Pr[\adv(\{\pk_i\}_{i\in[\pnum]}, \vec{v}_\pnum \allowbreak ) = \key] > 1/2 + \polysmall/2\}.$$ 
	A quick calculation yields $|\set| > 2^{Q(\lambda) - 1}\polysmall$. 
	Since $\{\sigma_i\}_{i \in [\log (m)]}$ and $\{r_i\}_{i \in [\log(m)]}$ are sampled uniformly at random, 
%	and if all $\sigma_{i \in [\log(m)]}$ are guessed correctly (i.e., all $\sigma_{i \in \mathcal{X}}$ is correct), 
	we have 
	$$\Pr[(\sigma_{i} = r_{i} \cdot \tempkey)_{i \in [\log (m)]}] = 1/m.$$ 
	Notice that the set $\{r_{\mathcal{X}}\}_{\mathcal{X} \subseteq [\log(m)]}$ are pairwise independent, the probability that the majority of bits is incorrect is 
	$$\Pr[(\tempkey_i^* \neq \tempkey_i) \mid (\sigma_j = \tempkey \cdot r_j)_{j \in [\log (m)]} \land (\tempkey \in S)] \leq 1/(m\polysmall^2),$$
	for each $i \in [Q(\lambda)]$ by Chebyshev bound. Putting all these together we have
	\[\begin{split}
		\Pr[b_{\lambda}(\{\pk_i\}_{i\in[\pnum]}) = \tempkey]
		\geq&  \Pr[\tempkey \in S] \cdot \Pr[b_{\lambda}(\{\pk_i\}_{i\in[\pnum]}) = \tempkey \mid \tempkey \in S]\\
		=& \polysmall/(2m) \cdot \Pr[\forall i \in [Q(\lambda)], (\tempkey^*_i = \tempkey_i) \mid \tempkey \in S, (\sigma_j = \tempkey \cdot r_j)_{j \in [\log (m)]}] \\
		=& \polysmall/(2m)(1-\Pr[\exists i \in [Q(\lambda)]\ \mbox{s.t.}\ \tempkey^*_i \neq \tempkey_i \mid \tempkey \in S, (\sigma_j = \tempkey \cdot r_j)_{j \in [\log (m)]}])\\
		\geq& \polysmall/(2m)(1-Q(\lambda)\Pr[\tempkey^*_i \neq \tempkey_i \mid (\sigma_j = \tempkey \cdot r_j)_{j \in [\log (m)]}, \tempkey \in S])\\
		\geq& \polysmall/(2m)(1-Q(\lambda)/(m\polysmall^2)) = \polysmall/(2m) - Q(\lambda)/(2m^2\polysmall).
	\end{split}\]
	Since $Q(\lambda) = \Theta(\log(\lambda))$ and $\epsilon = 1/\lambda^{o(1)}$, the probability that $b_{\lambda}$ succeeds is at least $\polysmall^3/(8\Theta(\log(\lambda))) = 1/\lambda^{o(1)}$.
	Since $\adv$ runs in time $\tilde{O}(T_2(\lambda))$ and $m = \Theta(\log \lambda) \lambda^{o(1)}$, $\{b_{\lambda}\}$ runs in time $\tilde{O}(\Theta(\log \lambda))\cdot m \cdot T_2(\lambda))=\tilde{O}(T_2(\lambda))$, completing the proof of \cref{thm:hcp}.
\end{proof}


%theorem of \epsilon/4 (with additional property)
%Let $\cc_1$ (respectively, $\cc_2$) be the set of all word-RAM families such that for each $\mathcal{F} = \{f_{\lambda}\}\in\cc_1$ and each $\lambda\in\N$, $f_{\lambda}$ runs in time $\tilde{O}(T_1(\lambda))$ (respectively, $\tilde{O}(T_2(\lambda))$) for any constant $\pnum \geq 2, k \geq 3$. Let $\epsilon =1/ \lambda^{o(1)}$ and $\tnike = (\gen, \share)$ be a $\cc_1$-$\pnum$-NIKE that satisfies $\cc_2$-$\frac{\epsilon}{4}$-weak security with the shared key size $ Q(\lambda) = \log(\lambda^{\Omega(1)})$, we construct $\tnike^* = (\gen^*, \share^*)$ satisfying $\cc_2$-$\epsilon$-security as in \cref{fig:con-search-nike}.
%\begin{theorem}
%\label{thm:hcp}
%	If $\tnike$ is a $\cc_1$-$\pnum$-NIKE with $\cc_2$-$\frac{\epsilon}{4}$-weak security, then $\tnike^*$ is a $\cc_1$-$\pnum$-NIKE with $\cc_2$-$\epsilon$-security.
%\end{theorem}
%\begin{figure}[h!]
%\begin{center}
%\framebox{
%\small
%  \begin{tabular}{l|l}
%    \begin{minipage}[t]{0.4\textwidth}
%    	\underline{$\gen^*(\bot)$:}
%	\item $(\pk, \sk) \getsr \gen(\bot)$
%%	\item $\vec{v} \getsr \{0,1\}^{\log\elllambda^3}$
%	\item $\vec{v} \getsr \{0,1\}^{Q(\lambda)}$
%    	\item Return $(\pk^* =(\pk, \vec{v}),\sk^* = \sk)$
%      \end{minipage}&  \begin{minipage}[t]{0.4\textwidth}
%        
%    \underline{$\share^*(i,\sk_i^*,(\pk_j^*)_{j\in[\pnum]})$:}
%    \item Parse $\pk_j^*=(\pk_j,\vec v_j)$ for all $j \in [\pnum]$
%    	\item $\tempkey = \share(i,\sk_i^*,(\pk_j)_{j\in[\pnum]}))$
%	%\item Let $\vec{v} = \vec{v}_\pnum$
%	\item Return $\key^* = \tempkey \cdot \vec{v}_\pnum$
%    \end{minipage} 
%      \end{tabular}
%      }
%\end{center}
%\caption{Definition of $\tnike^*=\{\gen^*,\share^*\}$ with $\cc_2$-$\epsilon$-security.}
%\label{fig:con-search-nike}
%\end{figure}
%\begin{proof}
%	Let $m =  2Q(\lambda)/\epsilon^2$, and $\advA = \{\adv\} \in \cc_2$ be any adversary breaking the $\cc_2$-$\epsilon$-security of $\tnike^*$.
%%	where $\cc_2$ runs in time $\tilde{\Omega}(\elllambda T(\lambda))$. 
%	We construct an adversary $\advB = \{b_{\lambda}\} \in \cc_2$ breaking the $\cc_2$-$\frac{\epsilon}{4}$-weak security of $\tnike$.
%	
%	Firstly, $b_{\lambda}$ chooses $(r_1, \cdots, r_{\log (m)}) \getsr \{0,1\}^{Q(\lambda)}$. Then, for every nonempty subset $\mathcal{X} \subseteq [\log (m)]$, $b_{\lambda}$ computes $r_\mathcal{X} = \oplus_{i \in \mathcal{X}} r_i$ and $\sigma_\mathcal{X} = \oplus_{i \in \mathcal{X}} \sigma_i$ for all $\sigma_1, \cdots, \sigma_{\log m} \in \bin$. 
%	For every $i \in [Q(\lambda)]$, $b_{\lambda}$ runs 
%	$$\tempkey_i^\mathcal{X} = \sigma_\mathcal{X} \oplus \adv(\{\pk_i\}_{i\in[\pnum]}, r_\mathcal{X} \oplus e^{i}),$$ 
%	where $e^{i}$ denotes the $Q(\lambda)$-bit string with $1$ in the $i$'th position and 0 everywhere else, and sets $\tempkey_i^*$ as the majority of ${\tempkey_i^\mathcal{X}}$. Finally, $b_{\lambda}$ outputs $\tempkey^* = \tempkey_1^*|| \cdots ||\tempkey_{m}^*$ if $\key^* \cdot r_i = \sigma_i$ for all $i \in [\log(m)]$.
%
%	Let 
%	$$\set = \{\tempkey \big{|} \Pr[\adv(\{\pk_i\}_{i\in[\pnum]}, \vec{v}_\pnum \allowbreak ) = \key] > 1/2 + \polysmall/2\}.$$ 
%	A quick calculation yields $|\set| > 2^{Q(\lambda) - 1}\polysmall$. 
%	Note that $\advb$ only returns $\key = \key^*$ when $\key^* \cdot r_i = \sigma_i$ for all $i \in [\log(m)]$, we have $\Pr[(\sigma_j = \tempkey \cdot r_j)_{j \in [\log (m)]}] = 1$, where $\key = \share(1, \sk_1, (\pk_j)_{j \in [\pnum]})$ and $(\pk_i, \sk_i) \getsr \gen(\bot)$ for all $i \in [\pnum]$.
%%	Since $\{\sigma_i\}_{i \in [\log (m)]}$ are sampled uniformly at random, 
%%%	and if all $\sigma_{i \in [\log(m)]}$ are guessed correctly (i.e., all $\sigma_{i \in \mathcal{X}}$ is correct), 
%%	we have 
%%	$$\Pr[(\sigma_{i} = r_{i} \cdot \tempkey)_{i \in [\log (m)]}] = 1/m.$$ 
%	Since the set $\{r_{\mathcal{X}}\}_{\mathcal{X} \subseteq [\log(m)]}$ are pairwise independent, the probability that the majority of bits is incorrect is 
%	$$\Pr[(\tempkey_i^* \neq \tempkey_i) \big{|} (\sigma_j = \tempkey \cdot r_j)_{j \in [\log (m)]} \land (\tempkey \in S)] \leq \frac{1}{m\polysmall^2},$$ 
%	for each $i \in [Q(\lambda)]$ by chebyshev bound. Putting all these together we have
%	\[\begin{split}
%		\Pr[b_{\lambda}(\{\pk_i\}_{i\in[\pnum]}) = \tempkey]
%		\geq&  \Pr[\tempkey \in S] \cdot \Pr[b_{\lambda}(\{\pk_i\}_{i\in[\pnum]}) = \tempkey \big{|} \tempkey \in S]\\
%		=& \frac{\polysmall}{2} \cdot \Pr[\forall i \in [Q(\lambda)], (\tempkey^*_i = \tempkey_i) \big{|} \tempkey \in S, (\sigma_j = \tempkey \cdot r_j)_{j \in [\log (m)]}] \\
%		=& \frac{\polysmall}{2}(1-\Pr[\exists i \in [Q(\lambda)]\ \mbox{s.t.}\ \tempkey^*_i \neq \tempkey_i \big{|} \tempkey \in S, (\sigma_j = \tempkey \cdot r_j)_{j \in [\log (m)]}])\\
%		\geq& \frac{\polysmall}{2}(1-Q(\lambda)\Pr[\tempkey^*_i \neq \tempkey_i \big{|} (\sigma_j = \tempkey \cdot r_j)_{j \in [\log (m)]}, \tempkey \in S])\\
%		\geq& \frac{\polysmall}{2}(1-\frac{Q(\lambda)}{m\polysmall^2}) = \frac{\polysmall}{2} - \frac{Q(\lambda)}{2m\polysmall}.
%	\end{split}\]
%	Since $m = 2Q(\lambda)/\epsilon^2$, the probability that $b_{\lambda}$ succeeds is at least $\frac{\polysmall}{4}$. 
%%	Let $T(\lambda) = \lambda^{\pnum+k}$. Since $\adv$ runs in time $\tilde{O}(T(\lambda))$ and $m = \log \lambda^{\Omega(1)} \lambda^{o(1)}$, $\{b_{\lambda}\}$ runs in time $\tilde{O}(\log(\lambda^{\Omega(1)})\cdot m \cdot T(\lambda))=\tilde{O}(T(\lambda))$, completing the proof of \cref{thm:hcp}.
%	Also note that $\adv$ runs in time $\tilde{O}(T_2(\lambda))$ and $m = \log \lambda^{\Omega(1)} \lambda^{o(1)}$, then $\{b_{\lambda}\}$ runs in time $\tilde{O}(\log(\lambda^{\Omega(1)})\cdot m^2 \cdot T_2(\lambda))=\tilde{O}(T_2(\lambda))$, completing the proof of \cref{thm:hcp}.
%\end{proof}
